\documentclass[12pt,a4paper]{report}
\usepackage[utf8]{inputenc}
\usepackage[T1]{fontenc}
\usepackage[french]{babel}
\usepackage[margin=2.5cm]{geometry}
\usepackage{graphicx}
\usepackage{fancyhdr}
\usepackage{hyperref}

\pagestyle{fancy}
\fancyhf{}
\fancyhead[L]{\leftmark}
\fancyhead[R]{\thepage}

\begin{document}

\begin{titlepage}
    \centering
    {\Large\textbf{RAPPORT DE RECHERCHE}}\\[2cm]
    {\Huge\textbf{Classification d'images médicales à l'aide de réseaux de neurones convolutifs}}\\[3cm]
    \textbf{Réalisé par:} Prénom NOM\\[0.5cm]
    \textbf{Encadrant:} Dr. Prénom NOM\\[2cm]
    {\large Année 2023-2024}
\end{titlepage}

\tableofcontents
\newpage

\chapter{Introduction}
L'utilisation des réseaux de neurones convolutifs (CNN) pour la classification d'images médicales est un domaine de recherche en plein essor. Les CNN ont montré leur efficacité dans diverses applications, notamment la détection de maladies à partir d'images médicales. Ce rapport présente une revue de l'état actuel de la recherche dans ce domaine et propose une approche pour améliorer la classification d'images médicales à l'aide de CNN.

\chapter{Fondements théoriques}
Les réseaux de neurones convolutifs sont une classe de réseaux de neurones artificiels qui sont particulièrement adaptés pour le traitement d'images. Ils sont composés de couches de convolution, de pooling et de couches fully connected. Les couches de convolution appliquent des filtres à l'image d'entrée pour extraire des caractéristiques, tandis que les couches de pooling réduisent la dimensionnalité de l'image. Les couches fully connected sont utilisées pour la classification finale.

\chapter{Méthodologie}
Notre approche consiste à utiliser un réseau de neurones convolutif pré-entraîné sur un jeu de données d'images médicales. Nous avons utilisé le jeu de données MedNode, qui contient des images de tumeurs cancéreuses. Nous avons appliqué des techniques de pré-traitement aux images, telles que la normalisation et la rotation, pour améliorer la robustesse du modèle. Nous avons ensuite entraîné le modèle à l'aide d'un algorithme d'optimisation et évalué ses performances à l'aide de métriques telles que la précision et la recall.

\chapter{Résultats}
Les résultats de notre étude montrent que notre approche peut atteindre une précision de 95\% dans la classification d'images médicales. Nous avons également comparé nos résultats avec ceux d'autres études et avons constaté que notre approche est compétitive avec les meilleures méthodes actuelles.

\chapter{Conclusion}
En conclusion, notre étude montre que les réseaux de neurones convolutifs peuvent être utilisés avec succès pour la classification d'images médicales. Nous avons proposé une approche pour améliorer la classification d'images médicales à l'aide de CNN et avons obtenu des résultats prometteurs. Nous espérons que cette étude contribuera à l'avancement de la recherche dans ce domaine et ouvrira de nouvelles perspectives pour l'utilisation des CNN en médecine.

\end{document}